\documentclass[12pt,a4paper]{article}
\usepackage[utf8]{inputenc}
\usepackage[vietnam]{babel}
\usepackage{amsmath}
\usepackage{amsfonts}
\usepackage{amssymb}
\usepackage{graphicx}
\usepackage{hyperref}
\usepackage{enumitem}
\usepackage[T5]{fontenc} % <- THÊM DÒNG NÀY ĐỂ HẾT LỖI MẤT CHỮ
\usepackage{lmodern}


\title{\textbf{Tài liệu Dự án: ShieldCall}}
\author{Hệ thống bảo vệ giọng nói thời gian thực}
\date{\today}

\begin{document}

\maketitle

\section{Tên Dự án}
\textbf{ShieldCall} (Hệ thống phát hiện lừa đảo tích hợp trí tuệ nhân tạo).

\section{Mục đích dự án}
Dự án được xây dựng nhằm mục đích cung cấp một giải pháp bảo mật âm thanh thời gian thực cho người dùng di động.
Cụ thể, hệ thống tập trung vào việc:
\begin{itemize}
    \item Phát hiện các giọng nói giả mạo (\textbf{Deepfake}) trong các cuộc gọi hoặc giao tiếp âm thanh trực tiếp.
    \item Cảnh báo người dùng ngay lập tức khi phát hiện có rủi ro lừa đảo công nghệ cao.
    \item Đảm bảo tính riêng tư bằng cách xử lý toàn bộ dữ liệu âm thanh trực tiếp trên thiết bị (\textbf{On-device processing}) và tự động hủy dữ liệu trên RAM sau khi phân tích.
\end{itemize}

\section{Chi tiết kỹ thuật}
Hệ thống được phát triển dựa trên kiến trúc Hybrid kết hợp giữa hiệu năng cao của C++ và tính linh hoạt của Android Native (Kotlin).

\subsection{Thông số âm thanh đầu vào}
\begin{itemize}
    \item \textbf{Sample Rate:} 16,000 Hz (Tiêu chuẩn cho các mô hình nhận dạng giọng nói).
    \item \textbf{Channel:} Mono (Kênh đơn).
    \item \textbf{Format:} 16-bit PCM (Dữ liệu thô để tối ưu hóa xử lý tín hiệu).
\end{itemize}

\subsection{Công nghệ sử dụng}
\begin{itemize}
    \item \textbf{Ngôn ngữ lập trình:} Kotlin (Giao diện và quản lý vòng đời ứng dụng), C++ (Xử lý tín hiệu và chạy mô hình AI).
    \item \textbf{Giao tiếp:} Java Native Interface (JNI) để truyền dữ liệu âm thanh từ tầng Kotlin xuống C++ với độ trễ cực thấp.
    \item \textbf{AI Engine:} TensorFlow Lite C++ API tích hợp trực tiếp vào nhân hệ thống.
    \item \textbf{Quản lý bộ nhớ:} Sử dụng cấu trúc dữ liệu \textbf{Ring Buffer} (Đệm vòng tròn) trên RAM để lưu trữ tạm thời 2-3 giây âm thanh mà không làm đầy bộ nhớ máy.
\end{itemize}

\section{Quy trình hoạt động}
Quy trình xử lý của ShieldCall diễn ra theo các bước nghiêm ngặt sau:

\begin{enumerate}
    \item \textbf{Giai đoạn Thu âm (Capture):} Ứng dụng sử dụng \texttt{AudioRecord} để bắt luồng âm thanh thực tế từ Microphone.
    \item \textbf{Giai đoạn Đệm RAM (Buffering):} Dữ liệu PCM được đẩy xuống tầng C++ JNI và nạp vào \texttt{Ring Buffer}. Tại đây, dữ liệu cũ sẽ bị ghi đè bởi dữ liệu mới sau mỗi chu kỳ (ví dụ 2 giây).
    \item \textbf{Xử lý tín hiệu số (DSP):} Trước khi đưa vào mô hình AI, âm thanh được lọc nhiễu và chuẩn hóa thông qua các thuật toán DSP.
    \item \textbf{Chạy AI Inference:} Engine TensorFlow Lite sẽ phân tích các đặc trưng tần số (Spectrogram) để xác định xem giọng nói có các dấu hiệu của AI tạo ra hay không.
    \item \textbf{Phản hồi (Feedback):} Kết quả (Score) được trả ngược về tầng Kotlin. Giao diện người dùng sẽ cập nhật trạng thái: \textbf{An toàn (Xanh)}, \textbf{Nghi vấn (Cam)} hoặc \textbf{Nguy hiểm (Đỏ)}.
    \item \textbf{Giải phóng (Flush):} Khi người dùng dừng bảo vệ, toàn bộ vùng đệm trên RAM sẽ bị xóa sạch hoàn toàn để đảm bảo an toàn thông tin.
\end{enumerate}

\section{Lộ trình phát triển chi tiết và Tiêu chí nghiệm thu}

\subsection{Giai đoạn 1: Thiết lập hạ tầng và Cấu hình nhân (Core Infrastructure)}
\begin{itemize}
    \item \textbf{Các bước thực hiện:}
    \begin{enumerate}
        \item Cấu hình \texttt{build.gradle} để hỗ trợ NDK và liên kết với \texttt{CMakeLists.txt}.
        \item Thiết lập cấu trúc thư mục JNI chuẩn: \texttt{cpp/include} (header), \texttt{cpp/libs} (thư viện TFLite), \texttt{cpp/src} (mã nguồn).
        \item Nhúng thư viện TensorFlow Lite JNI (\texttt{.so}) cho các kiến trúc CPU (arm64-v8a, x86\_64).
        \item Viết boilerplate cho \texttt{native-lib.cpp} để kiểm tra việc load thư viện thành công.
    \end{enumerate}
    \item \textbf{Kết quả dự kiến:} Dự án biên dịch không lỗi; Logcat hiển thị thông báo ``Library loaded successfully'' khi ứng dụng khởi chạy.
\end{itemize}

\subsection{Giai đoạn 2: Phát triển luồng thu âm và Quản lý bộ nhớ đệm (Audio \& Memory Management)}
\begin{itemize}
    \item \textbf{Các bước thực hiện:}
    \begin{enumerate}
        \item Xây dựng lớp \texttt{AudioStreamController} trong Kotlin sử dụng \texttt{AudioRecord} với chế độ \texttt{VOICE\_COMMUNICATION}.
        \item Thiết lập Coroutine (IO Dispatcher) để đọc luồng PCM 16-bit liên tục từ Microphone.
        \item Xây dựng lớp \texttt{RingBuffer} bằng C++ để quản lý vùng nhớ đệm vòng tròn (Circular Buffer) trên RAM.
        \item Viết hàm JNI \texttt{nativeWritePCM} để chuyển mảng \texttt{ShortArray} từ Kotlin xuống C++ một cách hiệu quả.
    \end{enumerate}
    \item \textbf{Kết quả dự kiến:} Luồng âm thanh được ghi vào RAM liên tục; không xảy ra hiện tượng tràn bộ nhớ hoặc rò rỉ bộ nhớ (memory leak).
\end{itemize}

\subsection{Giai đoạn 3: Tích hợp Trí tuệ nhân tạo và Xử lý tín hiệu (AI \& DSP Integration)}
\begin{itemize}
    \item \textbf{Các bước thực hiện:}
    \begin{enumerate}
        \item Xây dựng module \texttt{DSP.cpp} để chuẩn hóa âm thanh (Normalization) và cắt khung hình (Framing).
        \item Triển khai lớp \texttt{ModelRunner.cpp} sử dụng TensorFlow Lite Interpreter để nạp mô hình \texttt{.tflite} từ thư mục \texttt{assets}.
        \item Thiết lập cơ chế kích hoạt Inference định kỳ (ví dụ: mỗi 400ms dữ liệu âm thanh mới).
        \item Tối ưu hóa đa luồng tại tầng C++ để việc chạy AI không gây nghẽn luồng thu âm.
    \end{enumerate}
    \item \textbf{Kết quả dự kiến:} Hệ thống trả về điểm số xác suất (0.0 - 1.0) đại diện cho độ tin cậy của giọng nói sau mỗi 400ms.
\end{itemize}

\subsection{Giai đoạn 4: Xây dựng hệ thống Phản hồi và Cảnh báo UI (Feedback System)}
\begin{itemize}
    \item \textbf{Các bước thực hiện:}
    \begin{enumerate}
        \item Triển khai cơ chế JNI Callback để tầng C++ gọi ngược lại hàm \texttt{onNativeAIResult} trong Kotlin.
        \item Sử dụng \texttt{Dispatchers.Main} để cập nhật giao diện người dùng ngay lập tức từ kết quả AI.
        \item Xây dựng logic phân loại rủi ro: An toàn (<50\%), Nghi vấn (50-80\%), Nguy hiểm (>80\%).
        \item Cập nhật màu sắc (Xanh/Cam/Đỏ) và thông báo cảnh báo trên màn hình.
    \end{enumerate}
    \item \textbf{Kết quả dự kiến:} Người dùng thấy được tỷ lệ Deepfake thay đổi theo thời gian thực khi đang nói chuyện.
\end{itemize}

\subsection{Giai đoạn 5: Tối ưu hóa, Bảo mật và Đóng gói (Optimization \& Security)}
\begin{itemize}
    \item \textbf{Các bước thực hiện:}
    \begin{enumerate}
        \item Xây dựng hàm \texttt{nativeRAMFlush} để xóa sạch dữ liệu âm thanh thô trên RAM khi dừng.
        \item Xử lý vòng đời ứng dụng (\texttt{onDestroy}) để giải phóng tài nguyên Microphone và bộ nhớ AI.
        \item Tối ưu hóa CPU bằng cách điều chỉnh Thread Priority cho luồng thu âm.
        \item Kiểm tra và xử lý các trường hợp ngoại lệ (mất quyền truy cập Mic, file model lỗi).
    \end{enumerate}
    \item \textbf{Kết quả dự kiến:} Ứng dụng hoạt động mượt mà; dữ liệu âm thanh được bảo mật hoàn toàn; không để lại dấu vết dữ liệu trên thiết bị.
\end{itemize}

\end{document}
